
\phantomsection\label{proof:sga3-V-5-b-general}%
By \hyperref[par:sga3-V-5-a]{part~a)}, we may suppose that \(d_1\) is locally free of rank \(n\).
Let \((Y,p)\) be a cokernel of \((d_0,d_1)\) in the category of all
ringed spaces. The argument at the end of \hyperref[proof:sga3-V-4.1-c]{4(c)} shows that, in order to
prove \hyperref[thm:sga3-V-4.1-i]{Theorem~4.1(i)}, it is enough to show that \(Y\) is a scheme and
that \(p:X_0\to Y\) is a morphism of schemes. By
Lemma~\hyperref[sga3:exp5:lemma:1.2]{1.2}, this question is local on \(Y\).
Let \(y\in Y\) and choose \(x\in X_0\) with \(p(x)=y\). If \(x\) has a
saturated affine open neighborhood \(U\), then \hyperref[sec:sga3-V-4]{Section~4} shows that
\(p(U)\) is an affine open subset of \(Y\), and
\(p|_U:U\to p(U)\) is a morphism of schemes. Thus it remains to prove
that every \(x\in X_0\) has a saturated affine open neighborhood. We
proceed as follows; the argument is taken from SGA~1, VIII,
Corollary~7.6.

\SGAThreeDiagram[0.90\textwidth]
  {figures/exp5/p268_section5_affine_saturated_chain.png}
  {The nested open subsets used to construct a saturated affine
  neighborhood of \(x\).}
  {fig:sga3-exp5-p268-affine-saturated-chain}
  {LNM 151, p.~270}

By \hyperref[thm:sga3-V-4.1-b]{condition 4.1(b)}, there is an affine open subset \(V\subseteq X_0\)
which contains \(d_1(d_0^{-1}(x))\).%
Put \(F=X_0\setminus V\). Since \(d_1\) is integral,
\(d_1(d_0^{-1}(F))\) is closed, and
\[
  V'=X_0\setminus d_1(d_0^{-1}(F))
\]
is the largest saturated open subset contained in \(V\). The set
\(V'\) is a neighborhood of the finite set
\(d_1(d_0^{-1}(x))\). Hence there is a section \(f\) of the structure
sheaf of \(V\), vanishing on \(V\setminus V'\), such that
\(d_1(d_0^{-1}(x))\) is contained in the principal open subset
\(V_f\subseteq V\). We shall prove that the largest saturated open
subset \((V_f)'\) contained in \(V_f\) is affine; it will then be the
required neighborhood.

Put \(Z(f)=V'\setminus V_f\). Then \(d_0^{-1}(Z(f))\) is the set of
points of
\[
  d_0^{-1}(V')=d_1^{-1}(V')
\]
at which \(d_0^\ast(f)\) vanishes. The morphism induced by \(d_1\)
from this common inverse image to \(V'\) is locally free of rank
\(n\).%
Lemma~\hyperref[lem:sga3-V-4.1.1]{4.1.1} therefore shows that
\(d_1(d_0^{-1}(Z(f)))\) is precisely the zero locus of the norm \(N\)
of \(d_0^\ast(f)\) relative to \(d_1\). It follows that
\[
  (V_f)'=V'\setminus d_1(d_0^{-1}(Z(f)))
\]
is the set of points of \(V_f\) at which \(N\) does not vanish.
Consequently, \((V_f)'\) is affine.

This proves \hyperref[thm:sga3-V-4.1-i]{Theorem~4.1(i)}; assertions \hyperref[thm:sga3-V-4.1-ii]{(ii)}, \hyperref[thm:sga3-V-4.1-iii]{(iii)}, and the first part
of \hyperref[thm:sga3-V-4.1-iv]{(iv)} are then clear. It remains to prove the consequences stated
at the end of \hyperref[thm:sga3-V-4.1-iv]{(iv)}; see \hyperref[bib:sga3-V-Ray67a]{[Ray67a]}, Theorem~1(iii).

By hypothesis, the groupoid \(X_\ast\) comes from an equivalence
relation
\[
  i:R\longrightarrow X_0\times X_0,
\]
where \(i\) is an immersion. We have shown that
\(R\) is effective (see \SGARefLink{sga3:iv:definition:3:3:2}{Expos\'e IV, 3.3.2}), and that
\[
  p:X_0\longrightarrow Y=X_0/R
\]
is surjective and finite locally free; in particular, it is faithfully
flat and of finite presentation.
